\chapter{Related Work}\label{ch:related}
% =============================================================================

\section{Diffusion Planners and Policies}\label{sec:rel:diffusion}

Diffusion was first used to plan whole trajectories \parencite{janner2022planning} and then to learn
visuomotor policies \parencite{chi2023diffusion}. Motion planning diffusion constructs cost functions
for safe robot motion and biases the sampler towards them \parencite{carvalho2025motion}. \ac{DPCC}
combines a trajectory diffusion model with a model-based projection at every denoising step and with
constraint tightening \parencite{romer2025diffusion}. Its planner runs as \ac{MPC}: a one-dimensional
temporal U-Net generates a state--action plan over a fixed horizon, a batch of candidate plans is
drawn, the plan is made dynamically feasible against an Euler-integrated model of the plant, and one
action is executed before the next plan is generated \parencite[Sec.~5]{romer2025diffusion}. It is
both the codebase this thesis starts from and the baseline it is measured against.

\section{Flow Matching and Few-Step Policies}\label{sec:rel:fm}

Flow matching \parencite{lipman2023flow,lipman2024guide} and rectified flow \parencite{liu2023flow}
replace the denoising chain by an ordinary differential equation, and robot policies have adopted them
to act in few steps. FlowPolicy uses consistency flow matching on 3-D point clouds
\parencite{zhang2024flowpolicy}; ManiFlow combines flow matching with consistency training to generate
dexterous actions in one or two steps from visual, language and proprioceptive inputs
\parencite{yan2025maniflow}; Streaming Flow Policy treats the action trajectory itself as the flow
trajectory \parencite{jiang2025streaming}. VFP captures multimodality with a variational latent prior
\parencite{zhai2025vfp}, real-time chunking enables smooth asynchronous execution of flow policies
\parencite{black2025rtc}, $\pi_0$ is a vision-language-action flow model for general robot control
\parencite{black2024pi0}, and flow policies have also been trained with policy gradients
\parencite{yi2026fpg}. Most of these policies enforce no hard constraints on the generated
trajectory.

On the generative-modelling side, MeanFlow learns average velocities for one-step generation
\parencite{geng2025meanflow}, Improved MeanFlow recasts its objective as a loss on the instantaneous
velocity with a network that predicts the average velocity \parencite{geng2025imf}, and $\alpha$-Flow
unifies trajectory flow matching, shortcut models and MeanFlow in one family trained with a curriculum
\parencite{zhang2025alphaflow}.

\section{Accelerating Generative Planners}\label{sec:rel:accel}

Implicit samplers shorten the chain of a trained diffusion model \parencite{song2021denoising}, and
distillation compresses it further: One-Step Diffusion Policy distils a pre-trained diffusion policy
into a single-step action generator \parencite{wang2024onedp}. Both accelerate a model after it is
trained. Consistency models \parencite{song2023consistency} and MeanFlow \parencite{geng2025meanflow}
instead learn a few-step generator directly, and with flow matching the step budget remains a choice at
deployment.

\section{Constrained and Safe Generative Planning}\label{sec:rel:safety}

Constraints on generated trajectories are enforced by guidance, by projection, or by acting inside the
sampler. Guidance biases the sampler with the gradient of a cost \parencite{carvalho2025motion}.
Projection maps a sample onto the feasible set after sampling \parencite{power2023sampling} or at every
step \parencite{christopher2024constrained,romer2025diffusion}. Among the methods that act inside the
sampler, SafeFlowMatcher combines flow matching with control barrier functions
\parencite{yang2025safeflowmatcher}; Physics-Constrained Flow Matching enforces arbitrary nonlinear
constraints in pretrained flow-based models without retraining \parencite{utkarsh2025pcfm}; SAD-Flower
generates safe, admissible and dynamically consistent trajectories \parencite{huang2025sadflower};
UniConFlow incorporates equality and inequality constraints in one constrained flow-matching framework
\parencite{yang2025uniconflow}; XFlowMP is a task-conditioned generative motion planner built on
Schr\"odinger bridges \parencite{nguyen2025xflowmp}; and SafeFlowMPC couples flow matching with online
optimisation on a real seven-degree-of-freedom manipulator, generating joint-position trajectories
that a real-time solver projects onto the safe set inside the sampling loop
\parencite[Sec.~III]{oelerich2026safeflowmpc}.

HardFlow reformulates hard-constrained sampling of flow-matching models as trajectory optimisation and
projects the plan the sampler is predicted to reach \parencite{li2025hardflow}. Its derivation
decomposes the sampling problem the way \ac{MPC} decomposes a horizon, into one subproblem per
sampling step, and the resulting procedure is stated as an algorithm
\parencite[Alg.~1]{li2025hardflow}. It applies this optimisation only in the later sampling steps. At
a small step budget that convention leaves no step before the final one, and the method reduces to
projecting a finished sample --- a case that, to the author's knowledge, has not been characterised.

HardFlow is evaluated on the same obstacle-avoidance task as \ac{DPCC}: it follows the setup and the
constraint geometry of \parencite{romer2025diffusion}, counts projection after every sampling step ---
the projection of \ac{DPCC} --- among its baselines, and reports, beside the safety rate and the
computation time, the number of steps to reach the target on the trials that end without a collision
\parencite[Sec.~VII-A, Table~II]{li2025hardflow}. Its differences from this thesis therefore lie in the
control loop rather than in the task: it plans a horizon of sixteen steps and executes eight of them before replanning,
fixes the step budget at ten Euler steps, states its plant as a linear model fitted to the training
data, and solves the resulting program with an interior-point solver
\parencite[App.~C]{li2025hardflow}. Its planner is likewise a temporal U-Net
\parencite{janner2022planning}.

\section{Vision-Conditioned Imitation Learning}\label{sec:rel:visual}

D3IL provides manipulation tasks with diverse human demonstrations, including vision-based variants
\parencite{jia2024towards}, and X-IL, from the same group, explores the design space of
imitation-learning policies \parencite{jia2025xil}. Visuomotor diffusion policies condition the
denoising network on image features \parencite{chi2023diffusion}. Image-generation transformers use
patch tokens \parencite{peebles2022dit}, while policy transformers can represent action sequences. Flow-matching policies follow the same path: VITA flows
directly from visual representations to latent actions \parencite{gao2025vita}, and $\pi_0$ conditions
a flow model on images and language \parencite{black2024pi0}.

\section{Learning-Based Planning and Control for Aerial Vehicles}\label{sec:rel:aerial}

Imitation learning for aerial vehicles covers language-guided short-range flight, for which UAV-Flow
Colosseo provides a real-world benchmark \parencite{wang2025uavflow}, and trajectory planning: CGD
guides diffusion policies for \ac{UAV} trajectories with constraints never encountered in training
\parencite{kondo2024cgd}. Below the planner, geometric tracking control on the special Euclidean group
tracks position and attitude without the singularities of local coordinates. Its
position-controlled flight mode states the thrust magnitude and moment vector that track a position
reference, the attitude they imply, and the stability of the resulting closed loop
\parencite[Eqs.~(19)--(23), Prop.~3]{lee2010geometric}; the cascade from a commanded force to
individual rotor thrusts through a fixed allocation matrix is the structure of
\parencite{mellinger2011minimum}. These are the equations the controller of
\autoref{sec:method:geometric} implements. Predictive sampling offers derivative-free model
predictive control in simulation \parencite{howell2022predictive}. The quadrotor model used here is the Skydio X2 from MuJoCo
Menagerie \parencite{zakka2022menagerie}.

\section{Positioning of This Thesis}\label{sec:rel:positioning}

The methods above differ less in their generative model than in the loop around it: whether a plan is
executed once or inside a receding horizon, what the plan commands, where the constraints act, and
which backbone carries the comparison. \autoref{tab:related-loops} states these for the works this
thesis builds on and is measured against.

\begin{table}[htpb]
  \caption[Planner, loop and constraint mechanism in the closest related work]{The closest related
  work by the loop it runs in, not by its generative model. ``Executed'' is the number of planned
  steps applied before the next plan is generated. None of the sources states the parameter count of
  its network.}
  \label{tab:related-loops}
  \centering\small
  \begin{tabularx}{\linewidth}{lLLL}
    \toprule
    Work & Constraints & Loop and plan & Backbone and capacity \\
    \midrule
    Diffuser \parencite{janner2022planning}
      & cost guidance, no hard constraint
      & plan of states and actions, replanned
      & temporal U-Net, size not reported \\
    Diffusion Policy \parencite{chi2023diffusion}
      & none
      & receding horizon over a chunk of actions
      & U-Net or transformer, size not reported \\
    \ac{DPCC} \parencite{romer2025diffusion}
      & projection of every denoising iterate, tightened
      & receding horizon, horizon $8$, $1$ executed, Euler plant model
      & temporal U-Net, size not reported \\
    HardFlow \parencite{li2025hardflow}
      & projection of the predicted final sample, later steps only
      & receding horizon, horizon $16$, $8$ executed, fitted linear plant model
      & temporal U-Net, size not reported \\
    SafeFlowMPC \parencite{oelerich2026safeflowmpc}
      & projection onto the safe set inside the sampling loop
      & receding horizon over joint-position trajectories
      & flow-matching policy, size not reported \\
    This thesis
      & one feasible set, applied to the iterate or to the predicted endpoint
      & receding horizon, horizon $8$, $1$ executed, Euler plant model
      & matched temporal U-Net backbone; capacity in \autoref{sec:setup:protocol} \\
    \bottomrule
  \end{tabularx}
\end{table}

Fast generative sampling, hard constraint enforcement, visual conditioning and aerial embodiment have
each been studied, but in a receding-horizon controller the cost of fast sampling and the cost of
constraint enforcement interact. The comparison here measures that interaction within the same
planning framework.
This thesis places instantaneous-velocity matching and the two average-velocity objectives inside the
constrained controller
of \ac{DPCC} and compares its per-step projection with endpoint projection \parencite{li2025hardflow}
on the same feasible set, first on the benchmark of \ac{DPCC} with its setup unchanged. It then carries
the resulting controller to a vision-conditioned alignment task and to a quadrotor benchmark, both built
for this thesis, and adds the step-budget condition under which endpoint projection reduces to
projecting a finished sample.

% =============================================================================
